% =============================================================
% BODY (ES)
% =============================================================

\section{Resolución De La Actividad}
\vspace{2mm}

El paper seleccionado para esta actividad es \textit{Análisis de sentimientos en redes sociales y buscadores online. Enfoque predictivo financiero: un mapeo sistemático de la literatura}. El trabajo analiza cómo el análisis de sentimientos, aplicado sobre redes sociales, buscadores online, noticias financieras y otras fuentes digitales, puede complementar los datos financieros tradicionales para mejorar la comprensión y predicción del comportamiento del mercado.

\vspace{2mm}

El artículo utiliza un mapeo sistemático de la literatura como estrategia de investigación para relevar el estado actual de las técnicas de análisis de sentimientos aplicadas al ámbito financiero. A partir de los trabajos analizados, es posible estudiar el problema abordado, el enfoque de solución considerado, las principales decisiones de diseño que pueden identificarse y los atributos de calidad relacionados.

\vspace{2mm}

A continuación, se responden las cuatro preguntas propuestas por la actividad a partir de la lectura y el análisis del paper.


\subsection{¿Cuál Era El Problema?}
\vspace{2mm}

El problema central identificado en el paper es la dificultad de predecir con precisión las variaciones de los mercados financieros. Según el artículo, estas variaciones no dependen únicamente de datos históricos o indicadores económicos tradicionales, sino también de múltiples factores externos, como cambios económicos globales, noticias, comportamiento de los inversores y percepciones sociales difundidas mediante plataformas digitales.

\vspace{2mm}

El paper plantea que muchos operadores financieros consideran que el rendimiento histórico puede utilizarse para anticipar rendimientos futuros. Sin embargo, el trabajo amplía esta perspectiva al señalar que el sentimiento de los usuarios también cumple un papel importante. En este sentido, el problema no es solamente estadístico o financiero, sino también informacional: los modelos tradicionales pueden resultar incompletos si no consideran datos no estructurados provenientes de redes sociales, noticias, buscadores online o foros financieros.

\vspace{2mm}

Además, el artículo destaca que los consumidores e inversores son cada vez más difíciles de predecir. Esto vuelve necesario incorporar fuentes de información capaces de reflejar opiniones, emociones y tendencias en tiempo real.

\vspace{2mm}

Desde la perspectiva del diseño de software, esta problemática puede interpretarse como la necesidad de construir soluciones capaces de integrar información heterogénea, procesar datos textuales y transformarlos en información útil para complementar los modelos utilizados en la toma de decisiones financieras.


\subsection{¿Qué Solución Propone?}
\vspace{2mm}

El paper no presenta una herramienta de software implementada como solución concreta. El \textbf{mapeo sistemático de la literatura} constituye la estrategia de investigación utilizada por el autor para estudiar el problema y analizar el estado actual del área.

\vspace{2mm}

A partir de los trabajos relevados, el enfoque de solución identificado consiste en \textbf{integrar el análisis de sentimientos con modelos predictivos financieros}, complementando los datos financieros tradicionales con información no estructurada obtenida desde redes sociales, noticias, buscadores y otras plataformas digitales.

\vspace{2mm}

El análisis de sentimientos permite clasificar automáticamente los textos según su polaridad, por ejemplo, positiva, negativa o neutra. Esta información puede convertirse en una señal adicional para interpretar el comportamiento del mercado y utilizarse junto con datos financieros dentro de modelos predictivos.

\vspace{2mm}

El enfoque se apoya en la idea de que las redes sociales, las noticias web, las tendencias de Google y las discusiones en foros pueden reflejar percepciones colectivas acerca de empresas, activos o contextos económicos. La combinación de estas fuentes con modelos predictivos busca mejorar la toma de decisiones, anticipar movimientos del mercado y contribuir a una gestión de riesgos más oportuna.

\vspace{2mm}

El paper también señala que la capacidad predictiva puede mejorar cuando se combinan múltiples fuentes de información. Por lo tanto, el enfoque no depende de una única fuente de datos, sino que propone complementar los indicadores financieros tradicionales con información textual proveniente del entorno digital.


\subsection{¿Qué Decisiones De Diseño Aparecen?}
\vspace{2mm}

El paper no presenta el diseño detallado de una aplicación concreta ni define una arquitectura de software específica. Sin embargo, a partir de los enfoques analizados pueden identificarse distintas decisiones relevantes desde la perspectiva del diseño de una solución basada en análisis de sentimientos aplicado a la predicción financiera.

\vspace{2mm}

En primer lugar, aparece la decisión de \textbf{integrar múltiples fuentes de información}. El enfoque contempla datos financieros tradicionales junto con información proveniente de redes sociales, noticias web, tendencias de Google y discusiones en foros. Esta decisión busca enriquecer la información disponible y evitar que las predicciones dependan exclusivamente del comportamiento histórico de los mercados.

\vspace{2mm}

En segundo lugar, se identifica la utilización de \textbf{procesamiento de lenguaje natural y análisis de sentimientos}. La información proveniente de las fuentes digitales se presenta principalmente como texto no estructurado, por lo que debe procesarse para obtener información interpretable, como la polaridad positiva, negativa o neutra de las opiniones.

\vspace{2mm}

Desde una perspectiva de diseño, una solución de este tipo requeriría mecanismos capaces de obtener, procesar, clasificar y representar la información textual antes de incorporarla a los modelos predictivos. El paper no especifica la implementación concreta de estos mecanismos, por lo que este aspecto debe entenderse como una necesidad de diseño derivada del enfoque estudiado y no como una arquitectura definida por el autor.

\vspace{2mm}

Otra decisión relevante es la \textbf{integración de la información de sentimiento con modelos predictivos financieros}. El análisis de sentimientos no se plantea como una fuente de decisión independiente, sino como información complementaria que puede combinarse con indicadores y datos financieros para mejorar la interpretación de las tendencias del mercado.

\vspace{2mm}

También aparece la decisión de aprovechar \textbf{información cercana al tiempo real}. Las redes sociales, las noticias y otras plataformas digitales permiten observar rápidamente cambios de opinión y tendencias emergentes. Una solución basada en este enfoque debe contemplar la incorporación frecuente de nueva información para mantener la utilidad de los resultados dentro de un contexto financiero cambiante.

\vspace{2mm}

El paper también plantea la posibilidad de desarrollar \textbf{estrategias personalizadas y adaptativas}. Las recomendaciones pueden ajustarse de acuerdo con los perfiles de riesgo, las preferencias y los comportamientos de los inversores, por lo que el diseño debe contemplar diferentes usuarios y contextos de decisión.

\vspace{2mm}

Finalmente, el enfoque requiere considerar aspectos como la \textbf{robustez frente al ruido y la ambigüedad del lenguaje}, la \textbf{transparencia de los resultados} y la \textbf{protección de la información financiera}. El artículo reconoce que los modelos predictivos no son infalibles y que deben trabajar con información subjetiva, mercados volátiles y riesgos asociados al uso de datos financieros.


\subsection{¿Qué Atributos De Calidad Intenta Mejorar?}
\vspace{2mm}

El atributo de calidad más evidente asociado al enfoque es la \textbf{precisión}. El paper analiza cómo la incorporación de información proveniente de redes sociales, buscadores y noticias puede complementar los datos financieros tradicionales y contribuir a obtener predicciones más ajustadas sobre las tendencias del mercado.

\vspace{2mm}

También resulta relevante la \textbf{oportunidad de la información}. Las fuentes digitales pueden proporcionar información en tiempo real, permitiendo detectar cambios de opinión, tendencias emergentes o reacciones públicas antes de que se reflejen completamente en los indicadores financieros tradicionales. Esto puede favorecer decisiones más rápidas por parte de inversores, analistas y empresas.

\vspace{2mm}

Otro atributo importante es la \textbf{adaptabilidad}. El artículo plantea que el análisis de sentimientos puede utilizarse para personalizar estrategias financieras según perfiles de riesgo, preferencias y comportamientos de los inversores. Por lo tanto, una solución basada en este enfoque debe ser capaz de adaptar sus recomendaciones a diferentes usuarios y contextos.

\vspace{2mm}

La \textbf{robustez} también resulta necesaria. El lenguaje utilizado en redes sociales y noticias puede ser ambiguo, subjetivo o contener información irrelevante, mientras que los mercados financieros son altamente variables. Una solución de este tipo debe ser capaz de distinguir información relevante del ruido y evitar que datos poco confiables produzcan conclusiones débiles.

\vspace{2mm}

La \textbf{seguridad y la privacidad} constituyen otros aspectos relevantes. En sus conclusiones, el paper señala la necesidad de garantizar la privacidad y seguridad de los datos financieros. Esta consideración resulta especialmente importante cuando los sistemas trabajan con información relacionada con usuarios, inversores o instituciones.

\vspace{2mm}

Finalmente, el artículo destaca la importancia de la \textbf{transparencia y la comprensibilidad}. La complejidad de los modelos predictivos debe equilibrarse con la capacidad del usuario para comprender los resultados obtenidos. Por lo tanto, no resulta suficiente que una solución produzca una predicción: también debe presentar sus resultados de manera suficientemente clara para apoyar una toma de decisiones responsable.


\section{Conclusión}
\vspace{2mm}

En síntesis, el paper aborda la dificultad de anticipar el comportamiento de los mercados financieros en contextos caracterizados por múltiples variables y una elevada volatilidad. Para estudiar esta problemática, el autor utiliza un mapeo sistemático de la literatura y, a partir de los trabajos analizados, identifica como enfoque relevante la integración del análisis de sentimientos con modelos predictivos financieros.

\vspace{2mm}

El enfoque estudiado busca complementar los datos financieros tradicionales con información no estructurada proveniente de redes sociales, buscadores, noticias y otras plataformas digitales. De esta manera, el sentimiento expresado por usuarios e inversores puede convertirse en una fuente adicional de información para interpretar tendencias y apoyar la toma de decisiones.

\vspace{2mm}

Desde la perspectiva del diseño de software, pueden reconocerse decisiones vinculadas con la integración de múltiples fuentes de datos, el procesamiento de lenguaje natural, la incorporación de información de sentimiento en modelos predictivos, el uso de información cercana al tiempo real, la personalización y la consideración de riesgos relacionados con la calidad y la protección de los datos.

\vspace{2mm}

Los principales atributos de calidad asociados al enfoque son la precisión, la oportunidad de la información, la adaptabilidad, la robustez, la seguridad, la privacidad, la transparencia y la comprensibilidad de los resultados.

\vspace{2mm}

Es importante destacar que el paper no define una arquitectura ni un diseño detallado de una solución concreta. Por este motivo, las decisiones identificadas corresponden a características y necesidades de diseño que pueden derivarse de los enfoques estudiados, sin atribuir al artículo componentes, tecnologías o patrones que no hayan sido especificados.